\chapter{Three Iron Rules for Selling Time}

Wealth freedom has already been defined in strict terms: it is the state in which a person no longer has to sell their time in order to meet the necessities of life. That definition does not make the road shorter. It makes the road visible.

Before that milestone, most people still must sell time. The immediate question is sharper: while I still must sell my time, how should I sell it?

Time is scarce, but that word is still too weak. Time is non-renewable. At the end of a life, one might measure value efficiency as:

\[
\mathrm{life\ value\ efficiency}=\frac{\mathrm{value\ created}}{\mathrm{time\ spent}}.
\]

A job, therefore, should not be viewed casually as ``just a job.'' If the job consumes your time, it consumes part of your life. A person who must sell time should at least become a careful trader.

There are three iron rules for selling time on the road to wealth freedom:

\begin{enumerate}
\item Growth is the fundamental key.
\item Pay attention to value, not valuation.
\item Patience matters more than almost anything.
\end{enumerate}

\section*{Growth Is the Fundamental Key}

Success is not the end. Wealth freedom is not the end. A high exam score, a promotion, a profitable year, a public victory, a title, a number in a bank account: all of these can be milestones. None of them should be treated as the final station.

Growth is the permanent necessity.

One decision rule is enough to begin:

\begin{quote}
Will this choice help me accumulate more capability?
\end{quote}

This question is colder and more useful than asking whether a choice is comfortable, prestigious, stable, or immediately profitable. The point is not to despise any occupation. The point is to examine the accumulation effect. The useful test is not the name of the role. The useful test is whether the role forces, enables, and rewards the accumulation of new capability.

Every choice should be inspected from the angle of capability:

\begin{itemize}
\item What existing ability will this choice improve?
\item What new ability will this choice force me to acquire?
\item What future choices will this ability make possible?
\end{itemize}

Many people ask, ``Can I use my existing skills in this new choice?'' That sounds practical, but it often treats current ability as the boundary of the self. The better question is:

\begin{quote}
After this choice, what existing ability will be strengthened, and what new ability will I gain?
\end{quote}

\section*{Value Before Valuation}

Like a stock, a person has both value and valuation. Value is the real substance: capability, judgment, reliability, creativity, endurance, usefulness, and the ability to create results. Valuation is the market's current opinion: salary, title, applause, reputation, rank, or price.

The market rarely knows value exactly. Ray Dalio's advice is useful here: do not let your valuation become far higher than your real value.

If life contained only one trade, the highest possible valuation would seem ideal. But life contains many trades. In a repeated game, valuation detached from value becomes a liability.

If you focus on valuation, you will begin to do things that raise valuation instead of things that raise value. Over a long enough period, valuation is anchored by value. It can fluctuate, sometimes violently, but it cannot float forever without support.

The iron rule is simple: build value and let valuation follow. Do not ignore the market. Do not be naive about price. But never let the market's temporary opinion become more important than the substance that must support it.

\section*{Patience Above All}

Patience is to time trading what compounding is to investing. Its power may be quiet at first, but it eventually shows itself through accumulation.

The impatient person naturally favors final-station thinking. They want the matter finished. They want the result visible. Because they want speed, they do not accumulate well. Because they want visible proof, they care too much about valuation. Because they want relief, they avoid thinking deeply about value.

A common excuse is: ``I am impatient only because I do not like what I am doing now. When I finally do what I truly love, I will become focused, efficient, and disciplined.''

But not being able to do exactly what one loves is the ordinary condition of human life. Precisely in that condition, patience is needed. Patience is not proven when everything is easy, attractive, and naturally rewarding. It is proven when the current stage is awkward, slow, and necessary.

Patience includes patience with one's current self. Not indulgence, not laziness, not resignation, but enough steadiness to keep working while the results are still poor.

\section*{The Bottom Line}

Growth protects you from treating wealth freedom as retirement from life. Value protects you from becoming addicted to market opinion. Patience protects you long enough for growth and value to compound.

Together they form the bottom line for selling time:

\begin{itemize}
\item Choose the trade that increases capability.
\item Build real value before demanding a higher valuation.
\item Stay patient through the period when accumulation is invisible.
\end{itemize}

Before selling another year, month, week, or hour, ask:

\begin{quote}
Will this trade of time make me more capable, more valuable, and more able to endure the compounding process?
\end{quote}

If the answer is no, the price may be higher than it appears. If the answer is yes, even an ordinary piece of work may become part of the road to freedom.

\sourcenote{Material source: 2.md.}
